\chapter{Related Work}
Prior work such as Li & Chen (2025) laid foundations, but The baseline only considers static data tiering. This paper introduces a dynamic access-based lifecycle management layer.

In reference to Mitchell (2025), the integration highlights a distinct improvement metric.


Furthermore, this architectural pattern facilitates distributed state management that aligns with modern cloud-native principles. 
By effectively decoupling the data ingestion from processing, the system is less prone to sudden temporal spikes. 
The integration between highly available computing nodes ensures robust load distribution, effectively combating single points of failure.
This approach builds inherently on asynchronous event-driven paradigms. 


Moreover, bridging the gap identified in Li & Chen (2025) necessitates this novel perspective: The baseline only considers static data tiering. This paper introduces a dynamic access-based lifecycle management layer.

As noted in Campbell (2026), scalability remains paramount. The system design strictly adheres to this, as evidenced by the empirical measurements.

Prior work such as Li & Chen (2025) laid foundations, but The baseline only considers static data tiering. This paper introduces a dynamic access-based lifecycle management layer.


Furthermore, this architectural pattern facilitates distributed state management that aligns with modern cloud-native principles. 
By effectively decoupling the data ingestion from processing, the system is less prone to sudden temporal spikes. 
The integration between highly available computing nodes ensures robust load distribution, effectively combating single points of failure.
This approach builds inherently on asynchronous event-driven paradigms. 


Prior work such as Li & Chen (2025) laid foundations, but The baseline only considers static data tiering. This paper introduces a dynamic access-based lifecycle management layer.


Furthermore, this architectural pattern facilitates distributed state management that aligns with modern cloud-native principles. 
By effectively decoupling the data ingestion from processing, the system is less prone to sudden temporal spikes. 
The integration between highly available computing nodes ensures robust load distribution, effectively combating single points of failure.
This approach builds inherently on asynchronous event-driven paradigms. 


Prior work such as Li & Chen (2025) laid foundations, but The baseline only considers static data tiering. This paper introduces a dynamic access-based lifecycle management layer.

In reference to Mitchell (2025), the integration highlights a distinct improvement metric.


Furthermore, this architectural pattern facilitates distributed state management that aligns with modern cloud-native principles. 
By effectively decoupling the data ingestion from processing, the system is less prone to sudden temporal spikes. 
The integration between highly available computing nodes ensures robust load distribution, effectively combating single points of failure.
This approach builds inherently on asynchronous event-driven paradigms. 


Prior work such as Li & Chen (2025) laid foundations, but The baseline only considers static data tiering. This paper introduces a dynamic access-based lifecycle management layer.


Furthermore, this architectural pattern facilitates distributed state management that aligns with modern cloud-native principles. 
By effectively decoupling the data ingestion from processing, the system is less prone to sudden temporal spikes. 
The integration between highly available computing nodes ensures robust load distribution, effectively combating single points of failure.
This approach builds inherently on asynchronous event-driven paradigms. 


Prior work such as Li & Chen (2025) laid foundations, but The baseline only considers static data tiering. This paper introduces a dynamic access-based lifecycle management layer.


Furthermore, this architectural pattern facilitates distributed state management that aligns with modern cloud-native principles. 
By effectively decoupling the data ingestion from processing, the system is less prone to sudden temporal spikes. 
The integration between highly available computing nodes ensures robust load distribution, effectively combating single points of failure.
This approach builds inherently on asynchronous event-driven paradigms. 


Moreover, bridging the gap identified in Li & Chen (2025) necessitates this novel perspective: The baseline only considers static data tiering. This paper introduces a dynamic access-based lifecycle management layer.

Prior work such as Li & Chen (2025) laid foundations, but The baseline only considers static data tiering. This paper introduces a dynamic access-based lifecycle management layer.

In reference to Carter (2026), the integration highlights a distinct improvement metric.


Furthermore, this architectural pattern facilitates distributed state management that aligns with modern cloud-native principles. 
By effectively decoupling the data ingestion from processing, the system is less prone to sudden temporal spikes. 
The integration between highly available computing nodes ensures robust load distribution, effectively combating single points of failure.
This approach builds inherently on asynchronous event-driven paradigms. 


Prior work such as Li & Chen (2025) laid foundations, but The baseline only considers static data tiering. This paper introduces a dynamic access-based lifecycle management layer.


Furthermore, this architectural pattern facilitates distributed state management that aligns with modern cloud-native principles. 
By effectively decoupling the data ingestion from processing, the system is less prone to sudden temporal spikes. 
The integration between highly available computing nodes ensures robust load distribution, effectively combating single points of failure.
This approach builds inherently on asynchronous event-driven paradigms. 


As noted in Rivera (2025), scalability remains paramount. The system design strictly adheres to this, as evidenced by the empirical measurements.

Prior work such as Li & Chen (2025) laid foundations, but The baseline only considers static data tiering. This paper introduces a dynamic access-based lifecycle management layer.


Furthermore, this architectural pattern facilitates distributed state management that aligns with modern cloud-native principles. 
By effectively decoupling the data ingestion from processing, the system is less prone to sudden temporal spikes. 
The integration between highly available computing nodes ensures robust load distribution, effectively combating single points of failure.
This approach builds inherently on asynchronous event-driven paradigms. 


Prior work such as Li & Chen (2025) laid foundations, but The baseline only considers static data tiering. This paper introduces a dynamic access-based lifecycle management layer.

In reference to Rivera (2025), the integration highlights a distinct improvement metric.


Furthermore, this architectural pattern facilitates distributed state management that aligns with modern cloud-native principles. 
By effectively decoupling the data ingestion from processing, the system is less prone to sudden temporal spikes. 
The integration between highly available computing nodes ensures robust load distribution, effectively combating single points of failure.
This approach builds inherently on asynchronous event-driven paradigms. 
